\documentclass[11pt]{article}
\usepackage[margin=1in]{geometry}
\usepackage{amsmath,amssymb,amsthm}
\usepackage[hidelinks]{hyperref}
\newtheorem{theorem}{Theorem}
\newtheorem{lemma}{Lemma}
\title{Nonsingularity of the periodic random-walk system\\
\large A solution to MathDB 357004}
\author{Principia Math In Progress}
\date{18 September 2026}
\begin{document}
\maketitle
\begin{abstract}
We prove the nonsingularity conjecture of Gerv\'e and Grigutis for
every positive integer premium and every number of seasons, under
the source's net profit, positive zero-mass, and simple-root assumptions.
A null vector would produce generating functions continuous on the
closed unit disk. Their coefficients tend to zero and satisfy a
positive recurrence. An elementary maximum principle then forces
the null vector to vanish. Roots on the unit circle are included.
\end{abstract}

\section*{Task 1: the original problem and its full solution}
Let $X_1,\ldots,X_N$ be independent, nonnegative integer-valued
random variables, and let $\kappa,N\geq1$. We retain the net profit
condition used to construct system (16) in \cite{GG}:
\begin{equation}\label{eq:assumptions}
 \mu=\sum_{j=1}^N\mathbb E X_j<\kappa N,\qquad
 \mathbb P(X_1+\cdots+X_N=0)>0.
\end{equation}
Write $G_X(z)=\mathbb E z^X$. The source proves that
$G_{X_1+\cdots+X_N}(z)-z^{\kappa N}$ has exactly $\kappa N-1$
zeros other than $1$ in the closed unit disk, counted with
multiplicity \cite[Lemma 4]{GG}. In Conjecture 1 these zeros
$\alpha_1,\ldots,\alpha_{\kappa N-1}$ are assumed simple.
The conjecture states that the coefficient matrix of system (16)
is nonsingular.

We specify that matrix without its large block display. This also
fixes all cyclic indices. Set
\[
 Y_1=X_N,\quad Y_j=X_{j-1}\ (2\leq j\leq N),\qquad
 g_j(z)=G_{Y_j}(z),\quad p_{j,r}=\mathbb P(Y_j=r),
\]
and $F_j(r)=\sum_{q=0}^r p_{j,q}$. In particular $p_{j,0}>0$,
by independence and \eqref{eq:assumptions}. For a vector
$v=(v_{j,i})_{1\leq j\leq N,\ 0\leq i<\kappa}$, define
\begin{equation}\label{eq:B}
 B_j(z)=\sum_{i=0}^{\kappa-1}v_{j,i}
                 \sum_{r=i}^{\kappa-1}z^r F_j(r-i).
\end{equation}
The first $\kappa N-1$ rows of the source matrix $A$ act by
\begin{equation}\label{eq:row}
 (Av)_\ell=\sum_{j=1}^N
 \frac{\prod_{r=1}^{j-1}g_r(\alpha_\ell)}
      {\alpha_\ell^{\kappa(j-1)}}B_j(\alpha_\ell),
 \qquad 1\leq\ell<\kappa N.
\end{equation}
Its final row acts by
\begin{equation}\label{eq:last}
 (Av)_{\kappa N}=\sum_{j=1}^N B_j(1).
\end{equation}
Indeed,
$\sum_{r=i}^{\kappa-1}F_j(r-i)
=\sum_{q=0}^{\kappa-1-i}p_{j,q}(\kappa-i-q)$,
which is exactly the final row printed in system (16).
Equations \eqref{eq:row} and \eqref{eq:last} are respectively the
homogeneous forms of source equations (13) and (11).

\begin{theorem}\label{thm:main}
Under \eqref{eq:assumptions} and the simple-root hypothesis above,
the $\kappa N$ by $\kappa N$ matrix $A$ is nonsingular.
\end{theorem}

\subsection*{Two elementary uniqueness facts}
\begin{lemma}\label{lem:coefficients}
If $C$ is holomorphic in the open unit disk and continuous on its
closure, its Taylor coefficients $h(n)=[z^n]C(z)$ tend to zero.
\end{lemma}
\begin{proof}
Put $M=\max_{|z|\leq1}|C(z)|$. Cauchy's estimate gives $|h(n)|\leq M$.
Uniform continuity implies
$\varepsilon_\rho=\max_{|z|\leq1}|C(z)-C(\rho z)|\to0$ as
$\rho\to1^-$. Again by Cauchy's estimate,
$|(1-\rho^n)h(n)|\leq\varepsilon_\rho$, so
$|h(n)|\leq M\rho^n+\varepsilon_\rho$.
First let $n\to\infty$, then $\rho\to1^-$.
\end{proof}

\begin{lemma}\label{lem:maximum}
Let $h_1,\ldots,h_N$ be complex sequences on the nonnegative
integers, each tending to zero. Extend the indices periodically,
so $h_0=h_N$. Suppose $p_{j,r}\geq0$,
$\sum_{r\geq0}p_{j,r}=1$, $p_{j,0}>0$, and
\begin{equation}\label{eq:harmonic}
 h_{j-1}(u)=\sum_{r=0}^{u+\kappa}
                 p_{j,r}h_j(u+\kappa-r),\qquad u\geq0.
\end{equation}
Then all the sequences vanish identically.
\end{lemma}
\begin{proof}
If not, their common supremum in modulus is a positive number $M$.
It is attained at some $(j-1,u)$, since there are finitely many
sequences and each tends to zero. Taking absolute values in
\eqref{eq:harmonic} bounds the right-hand side by $M$.
Because $p_{j,0}>0$, equality is possible only if
$|h_j(u+\kappa)|=M$: otherwise that positive-weight term makes
the bound strictly smaller than $M$. Repeating this argument
propagates the modulus $M$ to arguments $u+q\kappa$ for every
$q\geq0$, contradicting convergence to zero.
\end{proof}

\subsection*{From a null vector to a decaying recurrence}
\begin{proof}[Proof of Theorem \ref{thm:main}]
Suppose $Av=0$, and form \eqref{eq:B}. Put
\[
 K(z)=\prod_{r=1}^N g_r(z)-z^{\kappa N}.
\]
Extend $g_j,B_j$ periodically in $j$, and, for $1\leq j\leq N$, set
\begin{equation}\label{eq:H}
 H_j(z)=\sum_{\ell=1}^N z^{\kappa(N-\ell)}
           \left(\prod_{r=1}^{\ell-1}g_{j+r}(z)\right)B_{j+\ell}(z).
\end{equation}
Direct cancellation of consecutive terms gives the algebraic identity
\begin{equation}\label{eq:cyclic}
 g_jH_j-z^\kappa H_{j-1}=K B_j,\qquad H_0=H_N.
\end{equation}
For $N=1$ this reads $(g_1-z^\kappa)B_1=KB_1$, so the
one-season case is included.

In particular
\[
 H_N(z)=\sum_{j=1}^N z^{\kappa(N-j)}
                  \left(\prod_{r=1}^{j-1}g_r(z)\right)B_j(z).
\]
The root equations \eqref{eq:row} imply $H_N(\alpha_\ell)=0$.
No $\alpha_\ell$ is zero, since $K(0)=\prod_jp_{j,0}>0$.
At any such root, \eqref{eq:cyclic} propagates the vanishing
backwards from $H_N$ to every $H_j$. At $1$ we have
\[
 H_j(1)=\sum_{r=1}^N B_r(1)=0
\]
by \eqref{eq:last}. Thus every $H_j$ vanishes at every zero
of $K$ in the closed unit disk.

Define $C_j=H_j/K$ away from these zeros. The interior
singularities are removable by the simple-root hypothesis.
We check the boundary carefully: finite mean in
\eqref{eq:assumptions} implies that each $g_j$ and $g_j'$
extends continuously to the closed disk, by absolute uniform
convergence of the derivative series. The same holds for $H_j,K$
and their first derivatives. Also $K'(1)=\mu-\kappa N<0$.
More generally, if $|\alpha|=1$ and $K(\alpha)=0$, equality in
the triangle inequality for the probability generating function
of $S=\sum_jY_j$, together with $\mathbb P(S=0)>0$, gives
\[
 \alpha^s=1\ \text{whenever }\mathbb P(S=s)>0,\qquad
 \alpha^{\kappa N}=1.
\]
Consequently
\begin{equation}\label{eq:boundary}
 K'(\alpha)=\alpha^{-1}(\mu-\kappa N)\ne0.
\end{equation}
For any function $f$ here vanishing at $\alpha$, integration
along the segment from $\alpha$ to $z$ gives
$f(z)/(z-\alpha)\to f'(\alpha)$ as $z\to\alpha$ within the
closed disk. Applying this to $H_j$ and $K$ proves that
$C_j$ extends continuously at every boundary root.
Hence $C_j$ is holomorphic in the disk and continuous on its closure.
Lemma \ref{lem:coefficients} implies
\[
 h_j(n):=[z^n]C_j(z)\longrightarrow0.
\]

Dividing \eqref{eq:cyclic} by $K$ and then continuing through
its removable zeros gives
\begin{equation}\label{eq:C}
 g_j(z)C_j(z)-z^\kappa C_{j-1}(z)=B_j(z).
\end{equation}
Since $\deg B_j<\kappa$, comparison of the coefficient of
$z^{u+\kappa}$ in \eqref{eq:C} gives exactly
\eqref{eq:harmonic}. Lemma \ref{lem:maximum} forces all
$h_j(n)$, hence all $C_j$ and all $B_j$, to vanish.
Finally the map from $(v_{j,0},\ldots,v_{j,\kappa-1})$ to
the coefficients of $B_j$ is triangular, with diagonal
$F_j(0)=p_{j,0}>0$. Therefore every $v_{j,i}=0$.
The square matrix $A$ has trivial kernel and is nonsingular.
\end{proof}

\paragraph{The explicitly stated three-season inequality.}
For $\kappa=1,N=3$, put
$u_\ell=G_{X_3}(\alpha_\ell)/\alpha_\ell$ and
$v_\ell=G_{X_3+X_1}(\alpha_\ell)/\alpha_\ell^2$.
Factoring the positive zero-masses from the three columns of
the source matrix (26), then subtracting its final row from
the first two, gives
\[
 \det A=
 \mathbb P(X_1=0)\mathbb P(X_2=0)\mathbb P(X_3=0)
 \bigl((u_1-1)(v_2-1)-(u_2-1)(v_1-1)\bigr).
\]
Theorem \ref{thm:main} makes this nonzero, which is precisely
the additional inequality requested in the conjecture.

\section*{Tasks 2--5: proposed follow-up directions}
These are further questions, not additional solved claims.
\begin{description}
\item[Task 2.] Remove the simple-root restriction by determining
the exact nonsingularity statement for the source's confluent
matrix, whose repeated root rows are replaced by derivative rows.
\item[Task 3.] When some $\mathbb P(X_j=0)$ vanish, characterize
which columns and root constraints must be removed to obtain a
square nonsingular system, including seasons whose minimum claim
exceeds $\kappa$.
\item[Task 4.] Find quantitative bounds for the smallest singular
value of the row-normalized system matrix in terms of the net
profit margin, zero-masses, root separation, and a specified
uniform moment bound.
\item[Task 5.] Generalize the argument to claim distributions
modulated by a finite-state Markov chain, giving explicit
irreducibility and boundary conditions that guarantee uniqueness
of the corresponding generating-function system.
\end{description}

\paragraph{Verification and status.}
The accompanying exact checks cover the cyclic numerator identities,
the matrix's normalization row, finite Taylor coefficient recurrences,
and boundary-root identities. They supplement the analytic proof
for arbitrary distributions. The published 35-page primary was read
in full; Conjecture 1 remains on printed page 335. The earlier
32-page arXiv source was checked at its definition and target
passages, not read in full. The dated literature audit records
bounded prior-claim checks. Verification is by the authoring agent
only, without independent peer review.

\begin{thebibliography}{9}
\raggedright
\bibitem{GG} S. Gerv\'e and A. Grigutis,
\emph{Distribution of shifted discrete random walk generated by
distinct random variables and applications in ruin theory},
Modern Stochastics: Theory and Applications \textbf{11}(3) (2024),
323--357.
\href{https://doi.org/10.15559/24-VMSTA249}{doi:10.15559/24-VMSTA249}.
Earlier version:
\href{https://arxiv.org/abs/2211.14629}{arXiv:2211.14629v1} (2022).
\end{thebibliography}
\end{document}
